\section{Related Work}
\label{sec:related-work}

\paragraph{Computational Counterpoint}
Music composition has been a popular domain for applying computational
techniques, and counterpoint is particularly well-studied because it
has clear rules.
Existing work represents counterpoint rules in a wide variety of ways,
including if-then statements~\cite{schottstaedt-cp}, fuzzy
logic~\cite{yilmaz-fuzzy}, answer set programming~\cite{boenn-asp,everardo2011armin},
and dependent and refinement types~\cite{szamozvancev-mezzo,
cong-musictools, cong-weighted}.
These approaches all guarantee the correctness of compositions, but
their implementation does not typically come with a user interface.

\paragraph{Mixed-Initiative Creativity Support Tools}

Our design of CounterChoice draws from a few closely-related genres
of software articulated in the literature:
\begin{itemize}
\item Procedural content generation, which emerges from digital games,
conceptualizes human creative activities as possibilities for
computational simulation (e.g. for the design of game levels and
environments); % TODO cite
\item Mixed-initiative
creative systems, which place procedural generation into a human
interaction loop,
conceptualizing human 
users and computational systems as engaged in a {\em collaboration}
where each can contribute ideas and content to the creative
artifact being developed; 
% (TODO cite)
\item Casual Creators~\cite{compton-casual}, which
focus on enabling the autotelic experience of creative activities rather than
task productivity, and center users with minimal prior experience in the creative
domain.
The philosophy of Casual Creators has inspired tools for creating
tile maps~\cite{carpenter-tile}, hybrid craft~\cite{graves-eloominate},
and avatars~\cite{otto-avatar}.
\end{itemize}

% Our tool's design attempts to follow this philosophy in that
% it allows the user to exploratively compose counterpoint without
% requiring knowledge in music theory or computer science. However,
% the degree to which it meets the criteria must ultimately be determined
% by observing its use by users outside its set of developers.

% TODO: 
% - describe how Counterchoice fits these frameworks
% - discuss related work for music composition

TODO: Music-specific tools for procedural generation/mixed-initiative
creativity/casual creators

TODO: logic programming/constraint solving-specific examples of the same


